%% toc: P3 40 40.6
\chapter{Almost-Sure Encoding and Infinite Spectra}
\label{sec:p03-w03-p03-069-infinite-spectra}

In the random graph \(G(n,n^{-\alpha})\) with \(\alpha\) irrational, every
first-order sentence \(A\) holds almost surely or almost never
(\cref{thm:p03-w03-p03-065-shelah-spencer}).  Letting \(\alpha\) vary gives a
\(\{0,1\}\)-valued function \(f_A\) of \(\alpha\), and the rationals at which
\(f_A\) jumps form the \emph{spectrum} of \(A\).  This chapter constructs a
single first-order sentence whose spectrum is infinite.  The construction
starts from a second-order property whose limiting behaviour is easy to read
off, namely that the largest common neighbourhood of three vertices has even
size, and replaces each second-order quantifier by first-order quantifiers
over auxiliary vertices.  The replacement is not an equivalence on all finite
graphs.  It is an \emph{almost sure encoding}: the second-order property and
the first-order sentence have the same truth value with probability tending
to one.

\paragraph{Inputs and conclusions.}
The inputs are an exponent \(\alpha\in(0,1)\), the random graph
\(G(n,p)\) with \(p=n^{-\alpha}\) (\cref{def:p03-w03-p03-065-gnp}), and
graph properties expressed either in first-order logic over the adjacency
relation \(\sim\) with equality, or with second-order quantifiers over binary
relations.  The conclusions are statements about
\(\lim_{n\to\infty}\Pr[G(n,n^{-\alpha})\entails A]\) for each \(\alpha\) in
explicit open intervals, and from these, rational points that belong to the
spectrum.  Every conclusion is asymptotic.  None of them decides a property of
a particular finite graph.  The first-order sentence and the second-order
property constructed below disagree on infinitely many finite graphs
(\cref{thm:p03-w03-p03-069-not-first-order}).

\paragraph{Selected treatment.}
The construction is that of Section~9, \emph{Infinite Spectra via Almost Sure
Encoding}, of Joel Spencer's \emph{Logic and Random Structures}, in the
30-page version hosted on Spencer's New York University page
\cite[pp.~16--18]{Spencer2007LogicRandomStructuresAuthor}.  That version is
not the printed chapter of the 2007 Springer volume \emph{Finite Model Theory
and Its Applications}, which has not been consulted; every locator below
refers to the hosted version, whose sections are numbered without a chapter
prefix.  The definitions of \(f_A\) and of the spectrum are in its Section~8
\cite[p.~15]{Spencer2007LogicRandomStructuresAuthor}.  The source gives the
construction with outline arguments.  It derives its key lemma from the
almost sure theory \(T_\alpha\), and it outlines a second, counting argument
which, it says, needs further work to become a full proof
\cite[p.~17]{Spencer2007LogicRandomStructuresAuthor}.  The proofs in the
annex are complete proofs written for this chapter.  The key lemma is proved
by a direct probabilistic argument along the lines of the second outline, so
that nothing here depends on the completeness of \(T_\alpha\).  Where the
source's numerical bookkeeping needs a restriction, the restriction is stated
and proved (\cref{def:p03-w03-p03-069-bookkeeping}).

\section{Spectra}
\label{sec:p03-w03-p03-069-spectra}

\begin{hbdefinition}{def:p03-w03-p03-069-spectrum}{The function \(f_A\) and the spectrum \(\operatorname{Sp}(A)\)}
Let \(A\) be a property of finite graphs.  For irrational \(\alpha\in(0,1)\)
put
\[
  f_A(\alpha)=\lim_{n\to\infty}\Pr\bigl[G(n,n^{-\alpha})\entails A\bigr]
\]
whenever the limit exists.  The domain is restricted to irrational
\(\alpha\) to avoid the rational exponents
\cite[p.~15]{Spencer2007LogicRandomStructuresAuthor}.  When \(A\) is a
first-order sentence, \cref{thm:p03-w03-p03-065-shelah-spencer} makes
\(f_A\) defined at every irrational \(\alpha\in(0,1)\), with values in
\(\{0,1\}\).

A rational \(q\in(0,1)\) is a \emph{point of discontinuity} of \(f_A\) if
there is no real number \(c\) such that \(f_A(\beta)\to c\) as \(\beta\to q\)
through irrational \(\beta\) in the domain of \(f_A\); equivalently, \(f_A\)
has no extension to the domain together with \(q\) that is continuous at
\(q\).  The \emph{spectrum} \(\operatorname{Sp}(A)\) is the set of rational
points of discontinuity of \(f_A\)
\cite[p.~15]{Spencer2007LogicRandomStructuresAuthor}.
\end{hbdefinition}

The source states the formulation with points of discontinuity; the
limit form above is the precise reading used in the proofs.  For a
first-order \(A\), the source shows that \(f_A\) is constant near each
irrational \(\alpha\), because \(A\) or \(\neg A\) follows from finitely many
axioms of \(T_\alpha\), each depending only on the signs of finitely many
numbers \(v-e\alpha\); hence discontinuities can occur only at rationals
\cite[p.~15]{Spencer2007LogicRandomStructuresAuthor}.  The constructions below
do not use that fact.  They exhibit, near each listed rational, two
intervals on which \(f_A\) takes different constant values, and this already
excludes a limit.  The source records that ``there exists a \(K_4\)'' has
spectrum \(\{2/3\}\) and ``there exists a \(K_5\)'' has spectrum \(\{1/2\}\)
(\cref{thm:p03-w03-p03-065-clique}), that a Boolean combination of the two has
spectrum \(\{2/3,1/2\}\), and that Spencer had once conjectured every
spectrum to be finite \cite[pp.~15--16]{Spencer2007LogicRandomStructuresAuthor}.
\Cref{thm:p03-w03-p03-069-infinite-spectrum} refutes that conjecture.

\section{Common neighbourhoods of triples}
\label{sec:p03-w03-p03-069-common-neighbourhoods}

\begin{hbdefinition}{def:p03-w03-p03-069-k3k}{\(K_{3,k}\), \(N(x_1,x_2,x_3)\) and the largest common neighbourhood}
Let \(k\geq1\).  A \emph{\(K_{3,k}\)} in a graph is a set of distinct
vertices \(x_1,x_2,x_3;y_1,\ldots,y_k\) such that every \(y_j\) is adjacent
to all three \(x_i\) \cite[p.~16]{Spencer2007LogicRandomStructuresAuthor}.
Adjacencies among the \(x\)'s or among the \(y\)'s are allowed and not
examined.  ``There exists a \(K_{3,k}\)'' is a first-order sentence.

For distinct vertices \(x_1,x_2,x_3\) of a graph \(G\), \(N(x_1,x_2,x_3)\)
is the set of common neighbours of \(x_1,x_2,x_3\)
\cite[p.~16]{Spencer2007LogicRandomStructuresAuthor}.  It never contains an
\(x_i\), since adjacency is irreflexive.  For a graph with at least three
vertices let
\[
  \nu(G)=\max\bigl\{\,|N(x_1,x_2,x_3)| : x_1,x_2,x_3\text{ distinct vertices of }G\,\bigr\}.
\]
The notation \(\nu(G)\) is introduced here; the source speaks of the maximal
size of \(|N(x_1,x_2,x_3)|\).  So \(\nu(G)\geq k\) exactly when \(G\)
contains a \(K_{3,k}\).

For \(k\geq1\) write
\[
  I_k=\Bigl(\tfrac13+\tfrac1{k+1},\;\tfrac13+\tfrac1k\Bigr)\cap(0,1).
\]
The intervals \(I_k\) are pairwise disjoint, and their union is
\((1/3,1)\) with the points \(1/3+1/k\), \(k\geq2\), removed.
\end{hbdefinition}

The triple \(x_1,x_2,x_3\) is read as distinct.  The source does not say so
explicitly, but its claim about the maximal size of \(|N(x_1,x_2,x_3)|\) is
false for repeated vertices (\cref{def:p03-w03-p03-069-distinctness}).

\begin{hblemma}{thm:p03-w03-p03-069-k3k-safe}{proved=annex:p03-w03-p03-069-k3k-safe}{\(K_{3,k}\) as a rooted graph}
Let \(k\geq1\) and let \(\alpha\in(0,1)\) be irrational.  The rooted graph
\((\emptyset,K_{3,k})\) has \(v=3+k\) nonroots and \(e=3k\) edges, and the
following are equivalent:
\((\emptyset,K_{3,k})\) is sparse; \((\emptyset,K_{3,k})\) is safe;
\(\alpha<\tfrac13+\tfrac1k\).
\end{hblemma}

Here sparse and safe are as in \cref{def:p03-w03-p03-065-dense-sparse}.  The
source records the equivalence of the first two conditions with
\(v-e\alpha>0\) \cite[p.~16]{Spencer2007LogicRandomStructuresAuthor}, and
the general criterion \cref{thm:p03-w03-p03-065-safe-criterion} then gives the
threshold at irrational exponents.  The next theorem proves the threshold
directly, for every edge probability function.

\begin{hbtheorem}{thm:p03-w03-p03-069-k3k-threshold}{proved=annex:p03-w03-p03-069-k3k-threshold}{Threshold for a \(K_{3,k}\)}
Let \(k\geq1\).  The function \(n^{-1/3-1/k}\) is a threshold function
(\cref{def:p03-w03-p03-065-threshold}) for the property ``there exists a
\(K_{3,k}\)'': for \(p=p(n)\in[0,1]\), the property holds almost never if
\(p\ll n^{-1/3-1/k}\) and almost surely if \(p\gg n^{-1/3-1/k}\)
\cite[p.~16]{Spencer2007LogicRandomStructuresAuthor}.
\end{hbtheorem}

\begin{hbcorollary}{thm:p03-w03-p03-069-max-common}{proved=annex:p03-w03-p03-069-max-common}{The largest common neighbourhood}
Let \(k\geq1\), \(\alpha\in I_k\) (not necessarily irrational) and
\(p=n^{-\alpha}\).  Then almost surely \(\nu(G(n,p))=k\).
\end{hbcorollary}

\begin{hbwarning}{def:p03-w03-p03-069-distinctness}{The three vertices must be distinct}
If the triple may repeat a vertex, then \(N(x,x,x)\) is the neighbourhood of
\(x\).  For \(\alpha\in(0,1)\) and \(p=n^{-\alpha}\), almost surely some
vertex of \(G(n,p)\) has at least \((n-1)p/2\) neighbours, and
\((n-1)p/2\to\infty\); so the largest common neighbourhood of a possibly
repeated triple tends to infinity, and
\cref{thm:p03-w03-p03-069-max-common} fails for it.  The proof is in the
annex item for \cref{thm:p03-w03-p03-069-max-common}.
\end{hbwarning}

\section{The second-order property}
\label{sec:p03-w03-p03-069-second-order}

\begin{hbdefinition}{def:p03-w03-p03-069-even-bigger}{\(\operatorname{EVEN}(S)\), \(\operatorname{BIGGER}(S,T)\) and \(A^*\)}
Let \(S\) and \(T\) be sets of vertices of a graph, given by first-order
formulas.  In second-order logic with a quantifier \(\exists_R\) over binary
relations \(R\) on the vertex set, the source writes
\[
  \operatorname{EVEN}(S):\ \exists_R\,
  \forall_x\neg R(x,x)\wedge\forall_{x,y}\bigl(R(x,y)\leftrightarrow R(y,x)\bigr)
  \wedge\forall_{x\in S}\exists!_{y\in S}R(x,y)
\]
\cite[p.~16]{Spencer2007LogicRandomStructuresAuthor}: there is an
irreflexive symmetric relation under which every element of \(S\) has exactly
one partner in \(S\).  The property \(\operatorname{BIGGER}(S,T)\) says that
there is an irreflexive symmetric binary relation \(R\) such that every
\(y\in T-S\) has some \(x\in S-T\) with \(R(y,x)\), and there are no distinct
\(y_1,y_2\in T-S\) and \(x\in S-T\) with \(R(y_1,x)\) and \(R(y_2,x)\)
\cite[p.~16]{Spencer2007LogicRandomStructuresAuthor}.  The property \(A^*\)
is
\[
  \begin{aligned}
  A^*:\ \exists_{x_1,x_2,x_3}&\operatorname{EVEN}[N(x_1,x_2,x_3)]\\
  &\wedge\forall_{z_1,z_2,z_3}\operatorname{BIGGER}[N(x_1,x_2,x_3),N(z_1,z_2,z_3)]
  \end{aligned}
\]
\cite[p.~16]{Spencer2007LogicRandomStructuresAuthor}, with both triples of
vertices read as ranging over distinct triples, for the reason given in
\cref{def:p03-w03-p03-069-distinctness}.
\end{hbdefinition}

\begin{hblemma}{thm:p03-w03-p03-069-even-bigger}{proved=annex:p03-w03-p03-069-even-bigger}{What the second-order formulas say}
For finite sets \(S,T\) of vertices of a graph,
\(\operatorname{EVEN}(S)\) holds exactly when \(|S|\) is even, and
\(\operatorname{BIGGER}(S,T)\) holds exactly when \(|T|\leq|S|\).
Consequently, for every graph \(G\) with at least three vertices,
\(G\entails A^*\) exactly when \(\nu(G)\) is even.
\end{hblemma}

\begin{hbtheorem}{thm:p03-w03-p03-069-astar-spectrum}{proved=annex:p03-w03-p03-069-astar-spectrum}{The limiting behaviour of \(A^*\)}
Let \(k\geq1\) and \(\alpha\in I_k\).  Then
\(\lim_{n\to\infty}\Pr[G(n,n^{-\alpha})\entails A^*]\) is \(1\) if \(k\) is
even and \(0\) if \(k\) is odd.  Hence
\[
  \operatorname{Sp}(A^*)\cap\bigl[\tfrac13,1\bigr)
  =\{\tfrac13\}\cup\bigl\{\tfrac13+\tfrac1k : k\geq2\bigr\}.
\]
\end{hbtheorem}

The source describes \(A^*\) as having all values \(1/3+1/k\) as spectral
points \cite[p.~16]{Spencer2007LogicRandomStructuresAuthor}; the theorem
proves this, and also that \(1/3\), where these points accumulate, is a
spectral point.  The behaviour of \(A^*\) for \(\alpha<1/3\) is not
analysed here.  The source also asserts that \(A^*\) cannot be written in the
first-order language \cite[p.~16]{Spencer2007LogicRandomStructuresAuthor}.

\begin{hbtheorem}{thm:p03-w03-p03-069-not-first-order}{proved=annex:p03-w03-p03-069-not-first-order}{\(A^*\) is not first-order}
No first-order sentence in the language \(\{\sim\}\) holds in exactly those
finite graphs that satisfy \(A^*\).  Consequently every first-order sentence
disagrees with \(A^*\) on infinitely many pairwise non-isomorphic finite
graphs.
\end{hbtheorem}

\section{The almost sure encoding}
\label{sec:p03-w03-p03-069-encoding}

\begin{hbdefinition}{def:p03-w03-p03-069-ru}{The predicate \(R_u(x,y)\)}
With \(u\) a variable symbol, let
\[
  R_u(x,y):=\exists_v\,[v\sim x\wedge v\sim y\wedge v\sim u],
\]
the first-order formula saying that \(u,x,y\) have a common neighbour
\cite[p.~17]{Spencer2007LogicRandomStructuresAuthor}.  For fixed \(u\) it
defines a symmetric binary relation in \(x,y\).  The idea is to replace a
second-order \(\exists_R\) by a first-order \(\exists_u\), and \(R\) by
\(R_u\) \cite[p.~17]{Spencer2007LogicRandomStructuresAuthor}.
\end{hbdefinition}

The source states its Representation Lemma as follows.  For every \(s\) and
every symmetric irreflexive \(R\) on \(\{1,\ldots,s\}\) holding for \(l\)
pairs with \(l<k/3\), the sentence
\[
  \forall_{x_1,\ldots,x_s}\exists_u\bigwedge_{1\leq i<j\leq s}
  \bigl(R_u(x_i,x_j)\leftrightarrow R(i,j)\bigr)
\]
belongs to \(T=T_\alpha\) for all \(\alpha\) with
\(\tfrac13+\tfrac1k>\alpha>\tfrac13+\tfrac1{k+1}\)
\cite[p.~17]{Spencer2007LogicRandomStructuresAuthor}.  The \(x_i\) are read
as distinct, as the matching application requires.  The lemma below proves
the almost sure form of the statement, with a hypothesis that does not refer
to \(k\).

\begin{hblemma}{thm:p03-w03-p03-069-representation}{proved=annex:p03-w03-p03-069-representation}{Representation lemma, almost sure form}
Let \(\tfrac13<\alpha<\tfrac12\) and \(p=n^{-\alpha}\).  Let \(s\geq2\) and
let \(R\) be a symmetric irreflexive relation on \(\{1,\ldots,s\}\) holding
for exactly \(l\) unordered pairs, where \(l(3\alpha-1)<1\).  Then almost
surely, for all distinct vertices \(x_1,\ldots,x_s\) of \(G(n,p)\), there is
a vertex \(u\) with
\(R_u(x_i,x_j)\leftrightarrow R(i,j)\) for all \(1\leq i<j\leq s\).
More precisely, there is \(c>0\), and with \(\delta=1-l(3\alpha-1)>0\) the
probability that the conclusion fails is at most
\(o(1)+n^se^{-cn^{\delta}}\).
\end{hblemma}

\begin{hblemma}{thm:p03-w03-p03-069-configurations}{proved=annex:p03-w03-p03-069-configurations}{The rooted graphs behind the representation}
Let \(\alpha\in(1/3,1/2)\) be irrational, with sparse, safe and rigid as in
\cref{def:p03-w03-p03-065-dense-sparse}.
\begin{enumerate}
\item With \(R\) as in \cref{thm:p03-w03-p03-069-representation}, let
\((S,H)\) have roots \(1,\ldots,s\), a nonroot \(u\), and for each pair
\(i<j\) with \(R(i,j)\) a nonroot \(v_{ij}\) adjacent to \(i\), \(j\) and
\(u\).  Then \((S,H)\) has \(v=1+l\) nonroots and \(e=3l\) edges, and it is
sparse exactly when \(l(3\alpha-1)<1\); in that case it is safe.
\item The rooted graph with three roots and one nonroot adjacent to all three
is rigid.
\end{enumerate}
\end{hblemma}

The source uses both facts in its derivation from \(T_\alpha\)
\cite[p.~17]{Spencer2007LogicRandomStructuresAuthor}.  Safety yields
witnesses \(u\) through the generic extension axiom, and rigidity of a common
neighbour of three vertices excludes unwanted pairs through
\(1\)-genericity; the axiom and genericity are defined in the source's
Section~5 \cite[p.~10]{Spencer2007LogicRandomStructuresAuthor}.  Within
\(I_k\), \(3\alpha-1<3/k\), so \(l<k/3\) gives \(l(3\alpha-1)<1\).  For
\(k\geq6\), where also \(I_k\subseteq(1/3,1/2)\), the hypothesis of
\cref{thm:p03-w03-p03-069-representation} is therefore implied by the
source's hypothesis.  For \(k\leq5\) the intervals \(I_k\) lie above \(1/2\),
and there the source's statement holds for \(k\leq3\) and fails for
\(k=4,5\), as the next result shows.  The proofs of this chapter use the
lemma only for \(k\geq6\).  A sentence that holds almost surely lies in
\(T_\alpha\) whenever \(T_\alpha\) is complete and all its members hold almost
surely, as the source asserts \cite[pp.~10--12]{Spencer2007LogicRandomStructuresAuthor}:
otherwise its negation would be in \(T_\alpha\) and would hold almost surely
too.  In that way the almost sure form gives the source's statement for
\(k\geq6\).

\begin{hbcorollary}{thm:p03-w03-p03-069-representation-fails}{proved=annex:p03-w03-p03-069-representation-fails}{The representation above \(1/2\)}
Let \(\tfrac12<\alpha<1\) and \(p=n^{-\alpha}\).
\begin{enumerate}
\item For every \(s\geq2\), almost surely every distinct
\(x_1,\ldots,x_s\) have a vertex \(u\) with \(\neg R_u(x_i,x_j)\) for all
\(i<j\); this is the case \(l=0\).
\item For \(s=2\) and \(R(1,2)\), so \(l=1\), almost surely there are
distinct \(x_1,x_2\) such that no vertex \(u\) satisfies \(R_u(x_1,x_2)\).
\end{enumerate}
Consequently, on \(I_k\) the sentence of the source's Representation Lemma
holds almost surely for every admissible \(R\) when \(k\leq3\), where
\(l<k/3\) forces \(l=0\), and holds almost never for \(l=1\) when
\(k\in\{4,5\}\), where \(l=1<k/3\).
\end{hbcorollary}

\begin{hbdefinition}{def:p03-w03-p03-069-encoded-formulas}{The encoded formulas and the sentence \(A\)}
Write \(\bar x=(x_1,x_2,x_3)\), \(\bar z=(z_1,z_2,z_3)\), and
\[
  \begin{aligned}
  \operatorname{Dist}(\bar x)&:=x_1\neq x_2\wedge x_1\neq x_3\wedge x_2\neq x_3,\\
  \operatorname{Adj}(y;\bar x)&:=y\sim x_1\wedge y\sim x_2\wedge y\sim x_3,
  \end{aligned}
\]
so that \(\operatorname{Adj}(y;\bar x)\) says \(y\in N(\bar x)\).
\begin{enumerate}
\item \(\operatorname{EVEN}^u(\bar x)\) is the first-order formula: there
exist \(u_1,u_2\) such that for every \(y\) with \(\operatorname{Adj}(y;\bar x)\)
there is exactly one \(y'\neq y\) with \(\operatorname{Adj}(y';\bar x)\) and
\(R_{u_1}(y,y')\vee R_{u_2}(y,y')\).  In words, there exist \(u_1,u_2\)
such that every common neighbour \(y\) of \(x_1,x_2,x_3\) has a unique other
common neighbour \(y'\) such that \(y,y',u_1\) or \(y,y',u_2\) have a common
neighbour \cite[p.~18]{Spencer2007LogicRandomStructuresAuthor}.
\item Let
\[
  R^{*}(y,w):=R_{u_1}(y,w)\vee R_{u_2}(y,w)\vee R_{u_3}(y,w)\vee R_{u_4}(y,w).
\]
Then \(\operatorname{BIGGER}^u(\bar x;\bar z)\) is: there exist
\(u_1,u_2,u_3,u_4\) such that
\begin{enumerate}
\item every \(y\) with \(\operatorname{Adj}(y;\bar z)\wedge\neg\operatorname{Adj}(y;\bar x)\)
has some \(w\) with \(\operatorname{Adj}(w;\bar x)\wedge\neg\operatorname{Adj}(w;\bar z)\)
and \(R^{*}(y,w)\); and
\item there are no \(y_1\neq y_2\) and \(w\) with
\(\operatorname{Adj}(y_m;\bar z)\wedge\neg\operatorname{Adj}(y_m;\bar x)\) for
\(m=1,2\), \(\operatorname{Adj}(w;\bar x)\wedge\neg\operatorname{Adj}(w;\bar z)\),
\(R^{*}(y_1,w)\) and \(R^{*}(y_2,w)\).
\end{enumerate}
This is \(\operatorname{BIGGER}(N(\bar x),N(\bar z))\) with \(\exists_R\)
replaced by \(\exists_{u_1,u_2,u_3,u_4}\) and \(R\) by
\(R_{u_1}\vee R_{u_2}\vee R_{u_3}\vee R_{u_4}\)
\cite[p.~18]{Spencer2007LogicRandomStructuresAuthor}.
\item The first-order sentence \(A\) is
\[
  A:=\exists\bar x\,\bigl[\operatorname{Dist}(\bar x)\wedge
  \operatorname{EVEN}^u(\bar x)\wedge
  \forall\bar z\,\bigl(\operatorname{Dist}(\bar z)\rightarrow
  \operatorname{BIGGER}^u(\bar x;\bar z)\bigr)\bigr].
\]
\end{enumerate}
The names \(\operatorname{EVEN}^u\), \(\operatorname{BIGGER}^u\),
\(\operatorname{Dist}\), \(\operatorname{Adj}\) and \(R^{*}\) are introduced
here; the source writes the formulas out in words.
\end{hbdefinition}

\begin{hbtheorem}{thm:p03-w03-p03-069-encoding}{proved=annex:p03-w03-p03-069-encoding}{The almost sure encoding}
\begin{enumerate}
\item In every graph and for all distinct \(\bar x\) and \(\bar z\):
\(\operatorname{EVEN}^u(\bar x)\) implies that \(|N(\bar x)|\) is even, and
\(\operatorname{BIGGER}^u(\bar x;\bar z)\) implies
\(|N(\bar z)|\leq|N(\bar x)|\).
\item Let \(k\geq6\), \(\alpha\in I_k\) and \(p=n^{-\alpha}\).  Then almost
surely the following hold in \(G(n,p)\) for all distinct \(\bar x\) and
\(\bar z\): \(\operatorname{EVEN}^u(\bar x)\) holds exactly when
\(|N(\bar x)|\) is even, and \(\operatorname{BIGGER}^u(\bar x;\bar z)\) holds
exactly when \(|N(\bar z)|\leq|N(\bar x)|\).  Consequently
\(\Pr[G(n,p)\entails A\leftrightarrow A^*]\to1\).
\end{enumerate}
\end{hbtheorem}

\begin{hbtheorem}{thm:p03-w03-p03-069-infinite-spectrum}{proved=annex:p03-w03-p03-069-infinite-spectrum}{A first-order sentence with infinite spectrum}
Let \(k\geq6\) and let \(\alpha\in I_k\) be irrational.  Then
\(f_A(\alpha)=1\) if \(k\) is even and \(f_A(\alpha)=0\) if \(k\) is odd.
Consequently
\[
  \operatorname{Sp}(A)\cap\bigl[\tfrac13,\tfrac12\bigr)
  =\{\tfrac13\}\cup\bigl\{\tfrac13+\tfrac1k : k\geq7\bigr\},
\]
and in particular the first-order sentence \(A\) has infinite spectrum
\cite[pp.~16--18]{Spencer2007LogicRandomStructuresAuthor}.
\end{hbtheorem}

\begin{hbmethod}{def:p03-w03-p03-069-method}{Building an almost sure encoding}
The construction follows a fixed pattern
\cite[pp.~16--18]{Spencer2007LogicRandomStructuresAuthor}.
\begin{enumerate}
\item Choose a second-order property \(P\) whose limiting probability is
known on explicit intervals of \(\alpha\); here
\(P=A^*\) and the intervals are the \(I_k\).
\item Bound the size of every set to which a second-order quantifier is
applied; here every \(N(\bar x)\) has at most \(k\) elements almost surely
(\cref{thm:p03-w03-p03-069-max-common}).
\item Replace each \(\exists_R\) by \(\exists_{u_1,\ldots,u_m}\) and \(R\)
by \(R_{u_1}\vee\cdots\vee R_{u_m}\), with \(m\) large enough that every
needed relation splits into \(m\) parts of \(l\) pairs each with
\(l(3\alpha-1)<1\).  Side condition: the representation lemma applies only
to such \(l\), and each part must be represented exactly on the relevant
finite set, not merely contain the wanted pairs.
\item Check the two directions separately.  The direction from the
first-order formula to the second-order one should hold in every graph; the
converse is where the representation lemma is used, and it holds only
almost surely.
\end{enumerate}
\end{hbmethod}

\begin{hbwarning}{def:p03-w03-p03-069-bookkeeping}{Where the counting needs \(k\geq6\)}
The source says that one \(R_u\) represents relations with at most \(k/3\)
pairs, so that two of them reach just under \(2k/3\) pairs, which covers a
matching on at most \(k\) points, and that four of them cover an injection
with \(k\) pairs \cite[pp.~17--18]{Spencer2007LogicRandomStructuresAuthor}.
In the proof of \cref{thm:p03-w03-p03-069-encoding} each part has at most
\(\lceil k/4\rceil\) pairs, and \(\lceil k/4\rceil\leq k/3\) holds for every
\(k\geq6\) but fails for \(k=1,2,5\).  Independently,
\cref{thm:p03-w03-p03-069-representation} is proved only for
\(\alpha<1/2\), and \(I_k\subseteq(1/2,1)\) for \(k\leq5\).  So the
argument determines \(f_A\) on \(I_k\) only for \(k\geq6\), that is for
\(\alpha<1/2\).  For \(\alpha\in I_k\) with \(k\leq5\), and for
\(\alpha<1/3\), the value of \(f_A\) is not determined here, and
\(1/2=1/3+1/6\) is not shown to be a spectral point of \(A\).
\end{hbwarning}

\begin{hbexample}{def:p03-w03-p03-069-pi-seven}{The jump at \(11/24\)}
Take \(\alpha=\pi/7=0.4487\ldots\), the source's sample exponent
\cite[p.~7]{Spencer2007LogicRandomStructuresAuthor}.  Since
\(\tfrac13+\tfrac19=0.4444\ldots<\pi/7<\tfrac{11}{24}=\tfrac13+\tfrac18\),
\(\alpha\in I_8\).
\begin{itemize}
\item By \cref{thm:p03-w03-p03-069-max-common}, almost surely the largest
common neighbourhood of three distinct vertices has exactly \(8\) elements.
\item A perfect matching of \(8\) points has \(4\) pairs, split as \(2+2\).
Each part has \(l=2\), and \(2(3\alpha-1)=0.692\ldots<1\).  The rooted graph
of \cref{thm:p03-w03-p03-069-configurations} for \(l=2\) has \(v=3\),
\(e=6\) and \(v-e\alpha=0.307\ldots>0\).
\item An injection between parts of two common neighbourhoods has at most
\(8\) pairs, split into four parts of at most \(2\).
\item Hence \(f_A(\pi/7)=f_{A^*}(\pi/7)=1\).
\end{itemize}
For irrational \(\alpha\) in \(I_7=(10/21,11/24)\) the largest common
neighbourhood has \(7\) elements almost surely, and \(f_A(\alpha)=0\).  So
\(f_A\) jumps at \(11/24\), and \(11/24\in\operatorname{Sp}(A)\).
\end{hbexample}

\section{Limits of the construction}
\label{sec:p03-w03-p03-069-limits}

\begin{hbwarning}{def:p03-w03-p03-069-only-almost-sure}{An encoding is not an equivalence}
\(A\) and \(A^*\) agree with probability tending to one in
\(G(n,n^{-\alpha})\) for \(\alpha\in I_k\), \(k\geq6\), and they disagree on
infinitely many finite graphs (\cref{thm:p03-w03-p03-069-not-first-order}).
The encoding depends on the probability measure: it is proved for
\(p=n^{-\alpha}\) in these intervals and for no other edge probability.  Only
one direction of the encoding, part~1 of
\cref{thm:p03-w03-p03-069-encoding}, holds in every graph.
\end{hbwarning}

\begin{hbwarning}{def:p03-w03-p03-069-hypothesis-limit}{The hypothesis of the representation lemma}
The condition \(l(3\alpha-1)<1\) is exactly the condition that the rooted
graph of \cref{thm:p03-w03-p03-069-configurations} is sparse.  When
\(l(3\alpha-1)>1\) that rooted graph is dense, the expected number of its
extensions over a fixed root tuple tends to \(0\), and the counting in the
proof of \cref{thm:p03-w03-p03-069-representation} gives no witness.  The
lemma is not claimed in that range.  The restriction \(\alpha<1/2\) is used
to make every pair of vertices have many common neighbours; above \(1/2\) even
a single designated pair cannot always be represented
(\cref{thm:p03-w03-p03-069-representation-fails}).  The restriction
\(\alpha>1/3\) is used to make every triple have only boundedly many common
neighbours.
\end{hbwarning}

The source closes the section with a general notion.  Given a property \(P\)
in a larger language \(L^+\), a sentence \(A\) of a smaller language \(L\)
is an almost sure encoding of \(P\) if the probability that \(P\) and \(A\)
differ in truth value tends to \(0\) as the size of the model tends to
infinity, for a fixed probability measure such as \(G(n,p(n))\)
\cite[p.~18]{Spencer2007LogicRandomStructuresAuthor}.  The source reports
that Hella, Kolaitis and Luosto call two languages almost everywhere
equivalent when each property in either has such an encoding in the other,
and that they ask, for \(G(n,1/2)\), whether monadic existential second-order
logic is almost everywhere equivalent to monadic universal second-order
logic.  They conjecture a negative answer; the source calls it an open
problem and remarks that negative results about the existence of almost sure
encodings seem hard to prove \cite[p.~18]{Spencer2007LogicRandomStructuresAuthor}.
This is an open question, not a theorem.

\section{Relationships}
\label{sec:p03-w03-p03-069-relationships}

The constant-probability zero--one law of \cref{sec:p03-zero-one-law} and
its countable limiting model, the random graph, are a different phenomenon.
There every first-order sentence has the same limit for all edge
probabilities \(p\in(0,1)\) held constant, and no function of a varying
exponent arises.  The descriptive-complexity presentation of
\cref{sec:p09-presentation-n-c-002} likewise concerns limiting theories and
their countable models, not the spectrum.  The ingredients used here are the
random graph and threshold notions of \cref{def:p03-w03-p03-065-gnp},
\cref{def:p03-w03-p03-065-almost-sure} and
\cref{def:p03-w03-p03-065-threshold}, the moment bounds
\cref{thm:p03-w03-p03-065-moments}, the dense--sparse classification
\cref{def:p03-w03-p03-065-dense-sparse}, and the game criterion for
first-order inexpressibility \cref{thm:p03-w03-p03-038-criterion}.  The
almost sure theory \(T_\alpha\) in which the source phrases its lemma is
described with the sparse-regime notions in
\cref{def:p03-w03-p03-067-sparse-notions}.  The source returns to almost sure
encoding in its Section~12, \emph{Nonconvergence via Almost Sure Encoding}, and its
Section~13 is titled \emph{No Almost Sure Representation of Evenness}
\cite[pp.~22 and 25]{Spencer2007LogicRandomStructuresAuthor}; those sections
are not presented here.

\paragraph{Sources.}
Joel Spencer, \emph{Logic and Random Structures}, 30-page version hosted at
\url{https://cs.nyu.edu/faculty/spencer/papers/ircsallsubm.pdf}: Section~8
for \(f_A\) and \(\operatorname{Sp}(A)\) (p.~15), Section~9 for the
construction (pp.~16--18), Sections~2, 3 and 5 for \(G(n,p)\), sparse and
safe rooted graphs and genericity (pp.~4, 6 and 10)
\cite{Spencer2007LogicRandomStructuresAuthor}.  The printed chapter of the
2007 Springer volume has not been consulted.
